\documentclass[bachelor]{zufe}
\usepackage{setspace} % 用于设置行间距
%首航缩进
\usepackage{indentfirst}
%各种下划线和url
\usepackage{hyperref}
\usepackage{ulem}
% 图表设置
\usepackage{graphicx}
\usepackage{caption}
\usepackage{appendix}
\usepackage{lipsum}
%数学相关包
\usepackage{algorithm}
\usepackage[noend]{algpseudocode} % 使用 algpseudocode 包提供的接口
% 自定义跨页算法环境
\makeatletter
\newenvironment{breakablealgorithm}
{% \begin{breakablealgorithm}
		\begin{center}
			\refstepcounter{algorithm}% New algorithm
			\hrule height.8pt depth0pt \kern2pt% \@fs@pre for \@fs@ruled
			\renewcommand{\caption}[2][\relax]{% Make a new \caption
				{\raggedright\textbf{\textbf{Algorithm}~\thealgorithm} ##2\par}%
				\ifx\relax##1\relax % #1 is \relax
				\addcontentsline{loa}{algorithm}{\protect\numberline{\thealgorithm}##2}%
				\else % #1 is not \relax
				\addcontentsline{loa}{algorithm}{\protect\numberline{\thealgorithm}##1}%
				\fi
				\kern2pt\hrule\kern2pt
			}
		}{% \end{breakablealgorithm}
		\kern2pt\hrule\relax %\@fs@post% for \@fs@ruled
	\end{center}
}
\makeatother
%====================Tiks ====================
\usepackage{tikz}
\usetikzlibrary{shapes.geometric, arrows.meta, positioning}
\tikzstyle{process} = [rectangle, minimum width=3cm, minimum height=1cm, text centered, draw=black, fill=blue!10]
\tikzstyle{arrow} = [thick, ->, >=stealth]
%====================Tiks ====================

\usepackage{amsmath} % 支持数学符号
\usepackage{amssymb}
\usepackage{algorithmicx}
%处理数学公式中和黑斜体
\usepackage{latexsym,bm}
%定义书写公式环境
\usepackage{cases}
%一些特殊符号例如：※√
\usepackage{pifont}

% 数学命令
\makeatletter
\newcommand\dif{%  % 微分符号
	\mathop{}\!%
	\ifthu@math@style@TeX
	d%
	\else
	\mathrm{d}%
	\fi
}
\makeatother
% 表格中支持跨行
\usepackage{multirow}

% 固定宽度的表格。
\usepackage{tabularx}

% 跨页表格
\usepackage{longtable}
%给字母加粗
\usepackage[T1]{fontenc}

\usepackage{amsthm}
% 量和单位
\usepackage{siunitx}

\newcommand{\formatmain}{
	\setlength{\parskip}{0em}
	\renewcommand{\baselinestretch}{1.53}
	\interfootnotelinepenalty=10000}
%三线表
\usepackage{booktabs}

%页眉页脚
\usepackage{fancyhdr}
%三线表
\usepackage{threeparttable}
% 中英文摘要和关键字
%这里标题的设置是黑体 三号 居中
%新罗马字体
\newcommand{\romann}{\CJKfontspec{TIMES .TTF}\selectfont}
\newtheorem{theorem}{定理}

%重命名lst名字
\renewcommand{\lstlistingname}{代码}
%参考文献生成规则
\usepackage[
backend=biber,
style=gb7714-2015,
gbalign=gb7714-2015,
gbnamefmt=lowercase,
gbpub=false,
doi=false,
url=false,
eprint=false,
isbn=false,
]{biblatex}

\hypersetup{
	colorlinks=true,
	linkcolor=black,
	citecolor=black,
	urlcolor=black
}


\addbibresource{Reference.bib}

% % 使用中文编译 + XeLatex
% \usepackage[fontset=ubuntu]{ctex}

% 基本信息--------------------------------------------------------------------------------
% 在这里填写你的论文中文题目
\newcommand{\thesisTitle}{海岛环境对武学宗师成长的影响机理}
% 在这里填写你的论文英文题目
\newcommand{\thesisTitleEN}{Influence mechanism of island environment on the growth of martial arts masters}

% 若有副标题，则运行下一行的代码，若无副标题，则将下一行注释掉（在\haveSub{}最前面添加 % 号）
\haveSub{}

% 在这里填写你的论文中文副标题（没有只需注释掉\haveSub{}即可）
\newcommand{\thesisSubTitle}{基于桃花岛武学流派的研究}
% 在这里填写你的论文英文副标题（没有只需注释掉\haveSub{}即可）
\newcommand{\thesisSubTitleEN}{Research Based on Taohua island martial arts school}

% 在这里填写你的相关信息
\newcommand{\deptName}{信息管理与人工智能学院}
\newcommand{\majorName}{人工智能}
%如果是生僻字例如:垚,则需要在使用\MyFont{}将生僻字括起来，这里用的是宋体，详细可见data文件夹下的zufesetup的第五行
\newcommand{\yourName}{\MyFont{王垚垚}}
\newcommand{\yourStudentID}{210110900xxx}
\newcommand{\mentorName}{吴XX}
\newcommand{\className}{人工智能1班}
\newcommand{\Today}{2024年4月}
% 基本信息--------------------------------------------------------------------------------

% 文档开始
\begin{document}
	
	% 封面，没有特殊情况不需要修改
	% 自定义下划线命令
\newcommand\dunderline[3][-1pt]{{%
		\setbox0=\hbox{#3}%
		\ooalign{\copy0\cr\rule[\dimexpr#1-#2\relax]{\wd0}{#2}}}}

% 生成封面
\begin{titlepage}
	\makeatletter
	%是否有校外指导老师，有的话就运行12行，否则运行15行，下面42行处的判断同理
	\@ifundefined{externalMentorName}{
		% 校内毕设封面顶部间距
		\vspace*{-20mm}
	}{
		% 校外毕设封面顶部间距
		\vspace*{-15mm}
	}
	\centering
	
	\includegraphics[width=3cm]{figures/logo.png}\\
	\includegraphics[width=5cm, height=1cm]{figures/name.png}
	\vspace*{10mm}
	\begin{center}
  		\zihao{1}\textbf{\ziju{0.12}\songti{本\hspace{5mm}科\hspace{5mm}生\hspace{5mm}毕\hspace{5mm}业\hspace{5mm}论\hspace{5mm}文（设计）}}
		
		\vspace{20mm}
		
		\xiaoer\textmd{\heiti{题目：}\thesisname}
		
		\vspace{5mm}
		%判断是否有副标题
		\ifx\subtitle\undefined\else
		%如果有副标题
		\begin{spacing}{1.2}
			\sanhao\selectfont{\textmd{\kaishu{-----}\subtitle}}
		\end{spacing}
		\fi
		\vspace{30mm}
		
		\flushleft
		\makeatletter
		\@ifundefined{externalMentorName}
		{
			\makeatother
			% 生成校内毕设封面字段
			\begin{spacing}{1.8}
				\hspace{27mm}\songti\zihao{3}\selectfont{学生姓名：\dunderline[-10pt]{1pt}{\makebox[78mm][c]{\studentname}}}
				
				\hspace{27mm}\songti\zihao{3}\selectfont{学\hspace{11mm}号：\dunderline[-10pt]{1pt}{\makebox[78mm][c]{\studentid}}}
				
				\hspace{27mm}\songti\zihao{3}\selectfont{指导教师：\dunderline[-10pt]{1pt}{\makebox[78mm][c]{\advisor}}}
				
				\hspace{27mm}\songti\zihao{3}\selectfont{所在学院：\dunderline[-10pt]{1pt}{\makebox[78mm][c]{\deptName}}}
				
				\hspace{27mm}\songti\zihao{3}\selectfont{专业名称：\dunderline[-10pt]{1pt}{\makebox[78mm][c]{\major}}}
				
				\hspace{27mm}\songti\zihao{3}\selectfont{班\hspace{11mm}级：\dunderline[-10pt]{1pt}{\makebox[78mm][c]{\className}}}
			\end{spacing}
		}
		{
			\makeatother
			% 生成校外毕设封面字段
			\begin{spacing}{1.8}
				\hspace{19.4mm}\songti\zihao{3}\selectfont{学\hspace{11mm}院：\dunderline[-10pt]{1pt}{\makebox[77.4mm][c]{\deptName}}}
				
				\hspace{19.4mm}\songti\zihao{3}\selectfont{专\hspace{11mm}业：\dunderline[-10pt]{1pt}{\makebox[77.4mm][c]{\major}}}
				
				\hspace{19.4mm}\songti\zihao{3}\selectfont{学生姓名：\dunderline[-10pt]{1pt}{\makebox[77.4mm][c]{\studentname}}}
				
				\hspace{19.4mm}\songti\zihao{3}\selectfont{学\hspace{11mm}号：\dunderline[-10pt]{1pt}{\makebox[77.4mm][c]{\studentid}}}
				
				\hspace{19.4mm}\songti\zihao{3}\selectfont{指导教师：\dunderline[-10pt]{1pt}{\makebox[77.4mm][c]{\advisor}}}
			\end{spacing}
		}

		
		\vspace{25mm}
		\centering
		\zihao{4}\ziju{0.5}{\thesisdate}
	\end{center}
\end{titlepage}
	
	% 前置页面定义
	\frontmatter
	% 原创性声明，没有特殊情况不需要修改
	% 原创性声明页
% 无特殊要求，不用修改

\fancypagestyle{originality}{
  % 页眉高度
  \setlength{\headheight}{10pt}

  % 页眉和页脚（页码）的格式设定
  \fancyhf{}
  \fancyhead[]{}

  % 页眉分割线稍微粗一些
  \renewcommand{\headrulewidth}{0pt}
}

\pagestyle{originality}
% \topskip=0pt

% % 圆形数字编号定义
% \newcommand{\circled}[2][]{\tikz[baseline=(char.base)]
%   {\node[shape = circle, draw, inner sep = 1pt]
%   (char) {\phantom{\ifblank{#1}{#2}{#1}}};
%   \node at (char.center) {\makebox[0pt][c]{#2}};}}
% \robustify{\circled}

% 设置行间距
\setlength{\parskip}{0.4em}
\renewcommand{\baselinestretch}{1.41}

% 顶部空白
\vspace*{-6mm}

% 原创性声明部分
\begin{center}
  \heiti\zihao{2}\textmd{声明及论文使用的授权}
\end{center}

\vspace{10mm}


% 本部分字号为小三
\zihao{-3}

本人郑重声明所呈交的论文是我个人在导师的指导下独立完成的。除了文中特别加以标注和致谢的地方外，论文中不包含其他人已经发表或撰写的研究成果。

\vspace{15mm}

\begin{flushright}
  论文作者签名：\hspace{75mm}年\hspace{8mm}月\hspace{8mm}日
\end{flushright}

\vspace{40mm}

% 使用授权声明部分

\zihao{-3}

本人同意浙江财经大学有关保留使用学位论文的规定，即：学校有权保留送交论文的复印件，允许论文被查阅和借阅；学校可以上网公布全部内容，可以采用影印、缩印或其他复制手段保存论文。

\vspace*{15mm}

\begin{flushright}
  论文作者签名：\hspace{75mm}年\hspace{8mm}月\hspace{8mm}日
\end{flushright}

\newpage

	
	% 摘要（中英文）：根据自身论文，修改摘要的Tex文件
	
% 下面是中文摘要 3-19行不用管
\setstretch{2}
\vspace*{2em}

\begin{center}
	\zihao{3}\textbf{ABSTRACT}
	\vspace{2mm}
	\ifx\subtitle\undefined\else
	%如果有副标题
	\begin{spacing}{1.2}
		\sanhao\selectfont{\textmd{\kaishu{-----}\subtitle}}
	\end{spacing}
	\fi
\end{center}
%题目后空一行
\vspace{1em}
\setstretch{1.5}

\setlength{\parskip}{0em}

% 英文摘要正文从这里开始--------------------------------------------------
As economic globalization continues to advance in depth, the translation of academic literature on economics has increasingly demonstrated its significance. This report selects and translates the 2015 Nobel Prize laureate’s academic literature into Chinese. Based on the research of the selected text and its Chinese translation by combining the case analysis with the comparative analysis, the report attempts to focus on what translation strategies are more effective in terms of realizing the ideational, interpersonal, and textual functions of the source text. The selected materials of this translation practice report do not have any existing translation versions for reference and the present author uses computer-aided technology to conduct a data analysis of the full original text before translation so as to learn about the features of the original.
\par
This report is guided by the three meta-functions of Halliday’s systemic functional linguistics. Through the analysis, it is concluded that annotation of terminologies, division of long sentences, and conversion of passive voice can achieve the ideational functions of the translation; addition of personal pronouns, literal translation of modal verbs and substitution of mood adjuncts can realize the interpersonal function; omission of connectives and conjunctive adjuncts, reproduction of demonstrative pronouns and adverbs, and reorganization of sentence order and paragraphs can achieve the textual function. The practicability of these strategies has been confirmed by concrete examples in this report. For the translation of academic literature on economics, comprehension and expression are pivotal. The translator should be familiar with the characteristics and style of the text, proficient in translation skills with the help of computer-aided tools, and flexible in dealing with rendering difficulties. This report is also intended to give some implications to both the translation practice of and theoretical application to the similar texts in the future.
\par

\vspace{1em}
\textbf{\heiti\xiaosi{Key words:}}
%这里写你的关键词------------------------------------------------------
academic literature on economics; three meta-functions of language; translation strategies

\newpage
\setstretch{2}
\vspace*{2em}
%如果没有英文标题，跳过
%下面是英文摘要 37-52行不用管
\begin{spacing}{0.95}
	\centering
	\sanhao\textbf{\heiti{摘要}}
	\vspace{2mm}
	\ifx\subtitle\undefined\else
	%如果有副标题
	\begin{spacing}{1.2}
		\sihao\selectfont{\textmd{\kaishu{-----}\rmfamily\textbf{\ensubtitle}}}
	\end{spacing}
	\fi
\end{spacing}
\vspace{1em}
\setstretch{1.5}
\setlength{\parskip}{0em}

% 中文摘要正文从这里开始--------------------------------------------------
随着经济全球化不断纵深推进，经济学类学术文献的翻译愈发显示出其重要性。本报告选择2015年诺贝尔经济学奖获得者的学术讲演稿进行汉译，基于对所选文本及其汉译文的研究，结合个案分析法与对比分析法，试图重点考察经济学类学术文献的翻译过程中采用何种策略可以在译文中更为有效地实现概念功能、人际功能、语篇功能等三大语言纯理功能。本翻译实践报告的所选材料无任何现有译文可供参考，笔者为了解原文特点，在译前借助计算机辅助技术对全文进行了数据分析。
\par
本报告以韩礼德系统功能语言学理论的三大语言纯理功能为指导，通过分析，结论如下：对术语加注、拆分长句和转换被动语态能够实现译文的概念功能；增译人称代词、直译情态动词和替换语气附加语可以实现译文的人际功能；省却连词和连接附加语、再现指示代词和指示副词，以及重组句序和段落能够实现译文的语篇功能，这些策略的实用性已由本报告的翻译实例得到印证。对于经济学类学术文献的翻译，理解能力与表达方式至关重要，译者当熟悉文体特征，精于翻译技巧并辅以机辅工具，灵活应对翻译难点。本报告对今后类似文本的翻译实践及理论应用也具有一定的启示意义。
\par

\vspace{1em}
\textbf{\rmfamily\xiaosi{关键词：}}
%这里写你的英文关键词----------------------------------------------------
经济学类学术文献；三大语言纯理功能；翻译策略



	%是否需要换页自己判断
	\newpage
	
	% 目录,自动生成，没有特殊情况不需要修改
	% 论文目录
% 没有特殊需要不用修改

% 调整目录行间距
\renewcommand{\baselinestretch}{1.35}
% 目录
\mainmatter
\addtocontents{toc}{\protect\thispagestyle{empty}}
\pagenumbering{gobble}
\tableofcontents
\clearpage
\pagenumbering{arabic}

	% 正文开始
	\mainmatter
	% 正文 22 磅的行距
	\setlength{\parskip}{0em}
	\renewcommand{\baselinestretch}{1.53}
	
	% 修复脚注出现跨页的问题
	\interfootnotelinepenalty=10000
	
	% 引言，根据自身论文，修改引言的Tex文件
	% 引言
\chapter{引言}
	\section{论文主要部分的写法}
	
	研究生学位论文撰写，除表达形式上需要符合一定的格式要求外，内容方面上也要遵循一些共性原则。
	
	通常研究生学位论文只能有一个主题（不能是几块工作拼凑在一起），该主题应针对某学科领域中的一个具体问题展开深入、系统的研究，并得出有价值的研究结论。
	学位论文的研究主题切忌过大，例如，“中国国有企业改制问题研究”这样的研究主题过大，因为“国企改制”涉及的问题范围太广，很难在一本研究生学位论文中完全研究透彻。
	
	
		
		\subsection{论文的语言及表述}
		
		除国际研究生外，学位论文一律须用汉语书写。
		学位论文应当用规范汉字进行撰写，除古汉语研究中涉及的古文字和参考文献中引用的外文文献之外，均采用简体汉字撰写。
		
		国际研究生一般应以中文或英文书写学位论文，格式要求同上。
		论文须用中文封面。
		
		研究生学位论文是学术作品，因此其表述要严谨简明，重点突出，专业常识应简写或不写，做到立论正确、数据可靠、说明透彻、推理严谨、文字凝练、层次分明，避免使用文学性质的或带感情色彩的非学术性语言。
		
		论文中如出现一个非通用性的新名词、新术语或新概念，需随即解释清楚。
		
		
		
		\subsection{论文题目的写法}
		
		论文题目应简明扼要地反映论文工作的主要内容，力求精炼、准确，切忌笼统。
		论文题目是对研究对象的准确、具体描述，一般要在一定程度上体现研究结论，因此，论文题目不仅应告诉读者这本论文研究了什么问题，更要告诉读者这个研究得出的结论。
		例如：“在事实与虚构之间：梅乐、卡彭特、沃尔夫的新闻观”就比“三个美国作家的新闻观研究”更专业、更准确。
		
		
		
		\subsection{摘要的写法}
		
		论文摘要是对论文研究内容的高度概括，应具有独立性和自含性，即应是 一篇简短但意义完整的文章。
		通过阅读论文摘要，读者应该能够对论文的研究 方法及结论有一个整体性的了解，因此摘要的写法应力求精确简明。
		论文摘要 应包括对问题及研究目的的描述、对使用的方法和研究过程进行的简要介绍、 对研究结论的高度凝练等，重点是结果和结论。
		
		论文摘要切忌写成全文的提纲，尤其要避免“第 1 章……；第 2 章……；……”这样的陈述方式。
		
		
		
		\subsection{引言的写法}
		
		一篇学位论文的引言大致包含如下几个部分：
		1、问题的提出；
		2、选题背 景及意义；
		3、文献综述；
		4、研究方法；
		5、论文结构安排。
		\begin{itemize}
			\item 问题的提出：要清晰地阐述所要研究的问题“是什么”。
			\footnote{选题时切记要有“问题意识”，不要选不是问题的问题来研究。}
			\item 选题背景及意义：论述清楚为什么选择这个题目来研究，即阐述该研究对学科发展的贡献、对国计民生的理论与现实意义等。
			\item 文献综述：对本研究主题范围内的文献进行详尽的综合述评，“述”的同时一定要有“评”，指出现有研究状态，仍存在哪些尚待解决的问题，讲出自己的研究有哪些探索性内容。
			\item 研究方法：讲清论文所使用的学术研究方法。
			\item 论文结构安排：介绍本论文的写作结构安排。
		\end{itemize}
		
		
		
		\subsection{正文的写法}
		
		本部分是论文作者的研究内容，不能将他人研究成果不加区分地掺和进来。
		已经在引言的文献综述部分讲过的内容，这里不需要再重复。
		各章之间要存在有机联系，符合逻辑顺序。
		
		
		
		\subsection{结论的写法}
		
		结论是对论文主要研究结果、论点的提炼与概括，应精炼、准确、完整，使读者看后能全面了解论文的意义、目的和工作内容。
		结论是最终的、总体的结论，不是正文各章小结的简单重复。
		结论应包括论文的核心观点，主要阐述作者的创造性工作及所取得的研究成果在本领域中的地位、作用和意义，交代研究工作的局限，提出未来工作的意见或建议。
		同时，要严格区分自己取得的成果与指导教师及他人的学术成果。
		
		在评价自己的研究工作成果时，要实事求是，除非有足够的证据表明自己的研究是“首次”、“领先”、“填补空白”的，否则应避免使用这些或类似词语。

	\newpage
	
	% 章节模板，成文后将其注释掉即可
	% !TeX root = ../zufe_report.tex
\setlength{\baselineskip}{20pt}
% 样例章节

% 章和引用示例-------------------------------------------------------------------------------------
\chapter{一级题目}

\section{二级题目}

正文······\cite{GB/T16159—1996} % 引用示例

\subsection{三级题目}

正文······\cite{Sobieski}
% 章和引用示例-------------------------------------------------------------------------------------


% 表示例----------------------------------------------------------------------------------------
\section{\textcolor{blue}{\underline{\underline{表-示例}}}}

\textcolor{blue}{{自动生成LaTeX表工具: \url{https://www.tablesgenerator.com/}}}

\begin{table}[htbp]
  \linespread{1.5}
  \zihao{5}
  \songti
  \centering
  \caption{物流的概念和范围}\label{物流的概念和范围}
  \begin{tabular}{|l|l|}
  \hline
  \multicolumn{1}{|c|}{本 质} & \multicolumn{1}{c|}{过  程}  \\ \hline
  途径或方法                     & 规划、实施、控制                   \\ \hline
  目标                        & 效率、成本效益                    \\ \hline
  活动或作业                     & 流动与储存                      \\ \hline
  处理对象                      & 原材料、在制品、产成品、相关信息           \\ \hline
  范围                        & 从原点（供应商）到终点（最终顾客）          \\ \hline
  目的或目标                     & 适应顾客的需求（产品、功能、数量、质量、时间、价格） \\ \hline
  \end{tabular}
\end{table}

 \begin{table}[htbp]
   \linespread{1.5}
   \zihao{5}
   \songti
   \centering
   \caption{统计表}\label{统计表}
   \begin{tabular}{|l|l|l|l|l|}
   \hline
   产品  & 产量    & 销量    & 产值   & 比重    \\ \hline
   手机  & 11000 & 10000 & 500  & 50\%  \\ \hline
   电视机 & 5500  & 5000  & 220  & 22\%  \\ \hline
   计算机 & 1100  & 1000  & 280  & 28\%  \\ \hline
   合计  & 17600 & 16000 & 1000 & 100\% \\ \hline
   \end{tabular}
 \end{table}

\begin{table}[htbp]
  \linespread{1.5}
  \zihao{5}
  \songti
  \centering
  \caption{统计表}\label{统计表}
  % Please add the following required packages to your document preamble:
  % \usepackage{multirow}
  \begin{tabular}{|l|l|l|l|l|}
  \hline
  年度                    & 产品  & 产量    & 销量    & 产值  \\ \hline
  \multirow{2}{*}{2004} & 手机  & 11000 & 10000 & 500 \\ \cline{2-5} 
                      & 计算机 & 1100  & 1000  & 280 \\ \hline
  \multirow{2}{*}{2005} & 手机  & 16000 & 13000 & 550 \\ \cline{2-5} 
                      & 计算机 & 2100  & 1500  & 320 \\ \hline
  \end{tabular}
\end{table}
% 表示例----------------------------------------------------------------------------------------

\section{\textcolor{blue}{\underline{\underline{表格-示例}}}}


表应具有自明性。为使表格简洁易读，尽可能采用三线表，如表~\ref{tab:three-line}。
三条线可以使用 booktabs 宏包提供的命令生成。

\begin{table}[htbp]
	\centering
	\caption{三线表示例}
	\begin{tabular}{ll}
		\toprule
		文件名          & 描述                         \\
		\midrule
		zufereport.dtx   & 模板的源文件，包括文档和注释 \\
		zufereport.cls   & 模板文件                     \\
		zufereport-*.bst & BibTeX 参考文献表样式文件    \\
		\bottomrule
	\end{tabular}
	\label{tab:three-line}
\end{table}

表格如果有附注，尤其是需要在表格中进行标注时，可以使用 threeparttable 宏包。
研究生要求使用英文小写字母 a、b、c……顺序编号，本科生使用圈码 ①、②、③……编号。

\begin{table}[htbp]
	\centering
	\begin{threeparttable}[c]
		\caption{带附注的表格示例}
		\label{tab:three-part-table}
		\begin{tabular}{ll}
			\toprule
			文件名                 & 描述                         \\
			\midrule
			zufereport.dtx\tnote{a} & 模板的源文件，包括文档和注释 \\
			zufereport.cls\tnote{b} & 模板文件                     \\
			zufereport-*.bst        & BibTeX 参考文献表样式文件    \\
			\bottomrule
		\end{tabular}
		\begin{tablenotes}
			\item [a] 可以通过 xelatex 编译生成模板的使用说明文档；
			使用 xetex 编译 \file{zufereport.ins} 时则会从 \file{.dtx} 中去除掉文档和注释，得到精简的 \file{.cls} 文件。
			\item [b] 更新模板时，一定要记得编译生成 \file{.cls} 文件，否则编译论文时载入的依然是旧版的模板。
		\end{tablenotes}
	\end{threeparttable}
\end{table}

如某个表需要转页接排，可以使用 longtable 宏包，需要在随后的各页上重复表的编号。
编号后跟表题（可省略）和“（续）”，置于表上方。续表均应重复表头。

\singlespacing
\begin{longtable}{cccc}
	\caption{跨页长表格的表题}
	\label{tab:longtable} \\
	\toprule
	表头 1 & 表头 2 & 表头 3 & 表头 4 \\
	\midrule
	\endfirsthead
	\caption*{续表~\thetable\quad 跨页长表格的表题} \\
	\toprule
	表头 1 & 表头 2 & 表头 3 & 表头 4 \\
	\midrule
	\endhead
	\bottomrule
	\endfoot
	Row 1  & & & \\
	Row 2  & & & \\
	Row 3  & & & \\
	Row 4  & & & \\
	Row 5  & & & \\
	Row 6  & & & \\
	Row 7  & & & \\
	Row 8  & & & \\
	Row 9  & & & \\
	Row 10 & & & \\
\end{longtable}

% 伪代码示例------------------------------------------------------------------------------------

\section{\textcolor{blue}{\underline{\underline{伪代码-示例}}}}

\textcolor{blue}{修改 algorithmic 之间的代码就可以实现论文伪代码，已经考虑了伪代码跨页问题}
\begin{breakablealgorithm}
        \caption{Calculate $y = x^n$}
        \begin{algorithmic}[1] %每行显示行号
            \Require $n \geq 0 \vee x \neq 0$ 
            \Ensure  $y = x^n$
            \State $y \gets 1$
            \If{$n < 0$}
            \State $ X \gets 1 / x$
            \State $N \gets -n$
            \Else
            \State $X \gets x$
            \State $N \gets n$
            \EndIf
            \While{ $N \neq 0$ }
            \If{ $N$ is even }
            \State $X \gets x \times x$
            \State $N \gets N / 2$
            \Else[$N$ is odd]
            \State $y \gets y \times X$
            \State $N \gets N - 1$
            \EndIf
            \EndWhile
     \end{algorithmic}
    \end{breakablealgorithm}
% 伪代码示例------------------------------------------------------------------------------------

% 代码块示例------------------------------------------------------------------------------------
\section{\textcolor{blue}{\underline{\underline{代码块-示例}}}}
\textcolor{blue}{只写了 C++, Python, Java 三种语言的格式}

% C++ 代码引入 + 文件引入
\lstinputlisting[
style = C++,
caption = {test.cpp},
label = {test.cpp}
]{./papperCode/test.cpp}

% Java 代码引入 + 行内引入
\begin{lstlisting}[
style = Java,
caption = {test.jar},
label = {test.jar}
]
public class HelloWorld {
    public static void main(String[] args){
        System.out.println("Hello World!");
    }
}
\end{lstlisting}

% Python 代码引入 + 文件引入
\lstinputlisting[
style = Python,
caption = {test.py},
label = {test.py}
]{./papperCode/test.py}
% 代码块示例------------------------------------------------------------------------------------
	\newpage
	
	% 在此插入章节
	% 第一章
	% !TeX root = ../zufe_report.tex
\setlength{\baselineskip}{20pt}
\chapter{第一章}
\section{论文主要部分的写法}

研究生学位论文撰写，除表达形式上需要符合一定的格式要求外，内容方面上也要遵循一些共性原则。

通常研究生学位论文只能有一个主题（不能是几块工作拼凑在一起），该主题应针对某学科领域中的一个具体问题展开深入、系统的研究，并得出有价值的研究结论。
学位论文的研究主题切忌过大，例如，“中国国有企业改制问题研究”这样的研究主题过大，因为“国企改制”涉及的问题范围太广，很难在一本研究生学位论文中完全研究透彻。



\subsection{论文的语言及表述}

除国际研究生外，学位论文一律须用汉语书写。
学位论文应当用规范汉字进行撰写，除古汉语研究中涉及的古文字和参考文献中引用的外文文献之外，均采用简体汉字撰写。

国际研究生一般应以中文或英文书写学位论文，格式要求同上。
论文须用中文封面。

研究生学位论文是学术作品，因此其表述要严谨简明，重点突出，专业常识应简写或不写，做到立论正确、数据可靠、说明透彻、推理严谨、文字凝练、层次分明，避免使用文学性质的或带感情色彩的非学术性语言。

论文中如出现一个非通用性的新名词、新术语或新概念，需随即解释清楚。



\subsection{论文题目的写法}

论文题目应简明扼要地反映论文工作的主要内容，力求精炼、准确，切忌笼统。
论文题目是对研究对象的准确、具体描述，一般要在一定程度上体现研究结论，因此，论文题目不仅应告诉读者这本论文研究了什么问题，更要告诉读者这个研究得出的结论。
例如：“在事实与虚构之间：梅乐、卡彭特、沃尔夫的新闻观”就比“三个美国作家的新闻观研究”更专业、更准确。



\subsection{摘要的写法}

论文摘要是对论文研究内容的高度概括，应具有独立性和自含性，即应是 一篇简短但意义完整的文章。
通过阅读论文摘要，读者应该能够对论文的研究 方法及结论有一个整体性的了解，因此摘要的写法应力求精确简明。
论文摘要 应包括对问题及研究目的的描述、对使用的方法和研究过程进行的简要介绍、 对研究结论的高度凝练等，重点是结果和结论。

论文摘要切忌写成全文的提纲，尤其要避免“第 1 章……；第 2 章……；……”这样的陈述方式。



\subsection{引言的写法}

一篇学位论文的引言大致包含如下几个部分：
1、问题的提出；
2、选题背 景及意义；
3、文献综述；
4、研究方法；
5、论文结构安排。
\begin{itemize}
	\item 问题的提出：要清晰地阐述所要研究的问题“是什么”。
	\footnote{选题时切记要有“问题意识”，不要选不是问题的问题来研究。}
	\item 选题背景及意义：论述清楚为什么选择这个题目来研究，即阐述该研究对学科发展的贡献、对国计民生的理论与现实意义等。
	\item 文献综述：对本研究主题范围内的文献进行详尽的综合述评，“述”的同时一定要有“评”，指出现有研究状态，仍存在哪些尚待解决的问题，讲出自己的研究有哪些探索性内容。
	\item 研究方法：讲清论文所使用的学术研究方法。
	\item 论文结构安排：介绍本论文的写作结构安排。
\end{itemize}



\subsection{正文的写法}

本部分是论文作者的研究内容，不能将他人研究成果不加区分地掺和进来。
已经在引言的文献综述部分讲过的内容，这里不需要再重复。
各章之间要存在有机联系，符合逻辑顺序。



\subsection{结论的写法}

结论是对论文主要研究结果、论点的提炼与概括，应精炼、准确、完整，使读者看后能全面了解论文的意义、目的和工作内容。
结论是最终的、总体的结论，不是正文各章小结的简单重复。
结论应包括论文的核心观点，主要阐述作者的创造性工作及所取得的研究成果在本领域中的地位、作用和意义，交代研究工作的局限，提出未来工作的意见或建议。
同时，要严格区分自己取得的成果与指导教师及他人的学术成果。

在评价自己的研究工作成果时，要实事求是，除非有足够的证据表明自己的研究是“首次”、“领先”、“填补空白”的，否则应避免使用这些或类似词语。

	\newpage
	
	% 第二章，第三章······
	% !TeX root = ../zufe_report.tex
\setlength{\baselineskip}{20pt}
\section{表格示例}
\subsection{表格}

表应具有自明性。为使表格简洁易读，尽可能采用三线表，如表~\ref{tab:three-line}。
三条线可以使用 booktabs 宏包提供的命令生成。

\begin{table}[htbp]
	\centering
	\caption{三线表示例}
	\begin{tabular}{ll}
		\toprule
		文件名          & 描述                         \\
		\midrule
		zufereport.dtx   & 模板的源文件，包括文档和注释 \\
		zufereport.cls   & 模板文件                     \\
		zufereport-*.bst & BibTeX 参考文献表样式文件    \\
		\bottomrule
	\end{tabular}
	\label{tab:three-line}
\end{table}

表格如果有附注，尤其是需要在表格中进行标注时，可以使用 threeparttable 宏包。
研究生要求使用英文小写字母 a、b、c……顺序编号，本科生使用圈码 ①、②、③……编号。

\begin{table}[htbp]
	\centering
	\begin{threeparttable}[c]
		\caption{带附注的表格示例}
		\label{tab:three-part-table}
		\begin{tabular}{ll}
			\toprule
			文件名                 & 描述                         \\
			\midrule
			zufereport.dtx\tnote{a} & 模板的源文件，包括文档和注释 \\
			zufereport.cls\tnote{b} & 模板文件                     \\
			zufereport-*.bst        & BibTeX 参考文献表样式文件    \\
			\bottomrule
		\end{tabular}
		\begin{tablenotes}
			\item [a] 可以通过 xelatex 编译生成模板的使用说明文档；
			使用 xetex 编译 \file{zufereport.ins} 时则会从 \file{.dtx} 中去除掉文档和注释，得到精简的 \file{.cls} 文件。
			\item [b] 更新模板时，一定要记得编译生成 \file{.cls} 文件，否则编译论文时载入的依然是旧版的模板。
		\end{tablenotes}
	\end{threeparttable}
\end{table}

如某个表需要转页接排，可以使用 \usepackage{longtable} 宏包，需要在随后的各页上重复表的编号。
编号后跟表题（可省略）和“（续）”，置于表上方。续表均应重复表头。

\singlespacing
\begin{longtable}{cccc}
	\caption{跨页长表格的表题}
	\label{tab:longtable} \\
	\toprule
	表头 1 & 表头 2 & 表头 3 & 表头 4 \\
	\midrule
	\endfirsthead
	\caption*{续表~\thetable\quad 跨页长表格的表题} \\
	\toprule
	表头 1 & 表头 2 & 表头 3 & 表头 4 \\
	\midrule
	\endhead
	\bottomrule
	\endfoot
	Row 1  & & & \\
	Row 2  & & & \\
	Row 3  & & & \\
	Row 4  & & & \\
	Row 5  & & & \\
	Row 6  & & & \\
	Row 7  & & & \\
	Row 8  & & & \\
	Row 9  & & & \\
	Row 10 & & & \\
\end{longtable}


\subsection{算法}

算法环境可以使用 algorithms 或者 algorithm2e 宏包。

\renewcommand{\algorithmicrequire}{\textbf{输入：}\unskip}
\renewcommand{\algorithmicensure}{\textbf{输出：}\unskip}

\begin{algorithm}
	\caption{Calculate $y = x^n$}
	\label{alg1}
	\small
	\begin{algorithmic}
		\REQUIRE $n \geq 0$
		\ENSURE $y = x^n$
		
		\STATE $y \leftarrow 1$
		\STATE $X \leftarrow x$
		\STATE $N \leftarrow n$
		
		\WHILE{$N \neq 0$}
		\IF{$N$ is even}
		\STATE $X \leftarrow X \times X$
		\STATE $N \leftarrow N / 2$
		\ELSE[$N$ is odd]
		\STATE $y \leftarrow y \times X$
		\STATE $N \leftarrow N - 1$
		\ENDIF
		\ENDWHILE
	\end{algorithmic}
\end{algorithm}

	\newpage
	
	% 结论：根据自身论文，修改结论的Tex文件
	%  结论
\unnumchapter{结~~~~论}
\renewcommand{\thechapter}{结论}

\ctexset{
  section/number = \arabic{section}
}
%从这里开始编写------------------------------------------------------------
% 结论部分尽量不使用 \subsection 二级标题，只使用 \section 一级标题

% 这里插入一个参考文献，仅作参考
本文结论……。\mycite{李成智2004飞行之梦}

\textcolor{black}{结论作为毕业设计（论文）正文的最后部分单独排写，但不加章号。结论是对整个论文主要结果的总结。在结论中应明确指出本研究的创新点，对其应用前景和社会、经济价值等加以预测和评价，并指出今后进一步在本研究方向进行研究工作的展望与设想。结论部分的撰写应简明扼要，突出创新性。阅后删除此段。}

\textcolor{black}{结论正文样式与文章正文相同：宋体、小四；行距：22 磅；间距段前段后均为 0 行。阅后删除此段。}
	\newpage
	
	% 参考文献，无特殊要求不需要修改
	% 添加参考文献请使用 BibTex 格式，添加至reference.bib 中，并在正文中使用 \cite{xxx}
	% 成文后将% 参考文献

% 参考文献开始
\unnumchapter{参考文献使用规则}
\renewcommand{\thechapter}{参考文献}

% 设置参考文献字号为 5 号
\renewcommand*{\bibfont}{\zihao{5}}
% 设置参考文献各个项目之间的垂直距离为 0
\setlength{\bibitemsep}{0ex}
\setlength{\bibnamesep}{0ex}
\setlength{\bibinitsep}{0ex}
% 设置单倍行距
\renewcommand{\baselinestretch}{1.2}
% 设置参考文献顺序标签 `[1]` 与文献内容 `作者. 文献标题...` 的间距
\setlength{\biblabelsep}{0.5mm}
% 设置参考文献后文缩进为 0（与 Word 模板保持一致）
\renewcommand{\itemcmd}{
  \addvspace{\bibitemsep} % 恢复 \bibitemsep 的作用
  \mkgbnumlabel{\printfield{labelnumber}}
  \hspace{\biblabelsep}}
% 删除默认的「参考文献 / Reference」标题，使用上面定义的 section 标题

\textcolor{blue}{参考文献书写规范}

\textcolor{blue}{参考国家标准《信息与文献参考文献著录规则》【GB/T 7714—2015】，参考文献书写规范如下：}

\textcolor{blue}{\textbf{1. 文献类型和标识代码}}

\textcolor{blue}{普通图书：M}\qquad\textcolor{blue}{会议录：C}\qquad\textcolor{blue}{汇编：G}\qquad\textcolor{blue}{报纸：N}

\textcolor{blue}{期刊：J}\qquad\textcolor{blue}{学位论文：D}\qquad\textcolor{blue}{报告：R}\qquad\textcolor{blue}{标准：S}

\textcolor{blue}{专利：P}\qquad\textcolor{blue}{数据库：DB}\qquad\textcolor{blue}{计算机程序：CP}\qquad\textcolor{blue}{电子公告：EB}

\textcolor{blue}{档案：A}\qquad\textcolor{blue}{舆图：CM}\qquad\textcolor{blue}{数据集：DS}\qquad\textcolor{blue}{其他：Z}

\textcolor{blue}{\textbf{2. 不同类别文献书写规范要求}}

\textcolor{blue}{\textbf{期刊}}

\noindent\textcolor{blue}{[序号]主要责任者. 文献题名[J]. 刊名, 出版年份, 卷号(期号): 起止页码. }


\textcolor{blue}{\textbf{普通图书}}

\noindent\textcolor{blue}{[序号]主要责任者. 文献题名[M]. 出版地: 出版者, 出版年. 起止页码. }
\mycite{Raymer1992Aircraft}


\textcolor{blue}{\textbf{会议论文集}}

\noindent\textcolor{blue}{[序号]析出责任者. 析出题名[A]. 见(英文用In): 主编. 论文集名[C]. (供选择项: 会议名, 会址, 开会年)出版地: 出版者, 出版年. 起止页码. }
\mycite{sunpinyi}


\textcolor{blue}{\textbf{专著中析出的文献}}

\noindent\textcolor{blue}{[序号]析出责任者. 析出题名[A]. 见(英文用In): 专著责任者. 书名[M]. 出版地: 出版者, 出版年.起止页码. }
\mycite{luoyun}


\textcolor{blue}{\textbf{学位论文}}

\noindent\textcolor{blue}{[序号]主要责任者. 文献题名[D]. 保存地: 保存单位, 年份. }
\mycite{zhanghesheng,Sobieski}


\textcolor{blue}{\textbf{报告}}

\noindent\textcolor{blue}{[序号]主要责任者. 文献题名[R]. 报告地: 报告会主办单位, 年份. }
\mycite{fengxiqiao,Sobieszczanski}


\textcolor{blue}{\textbf{专利文献}}

\noindent\textcolor{blue}{[序号]专利所有者. 专利题名[P]. 专利国别: 专利号, 发布日期. }
\mycite{jiangxizhou}


\textcolor{blue}{\textbf{国际、国家标准}}

\noindent\textcolor{blue}{[序号]标准代号. 标准名称[S]. 出版地: 出版者, 出版年. }
\mycite{GB/T16159—1996}


\textcolor{blue}{\textbf{报纸文章}}

\noindent\textcolor{blue}{[序号]主要责任者. 文献题名[N]. 报纸名, 出版年, 月(日): 版次. }
\mycite{xiexide}


\textcolor{blue}{\textbf{电子文献}}

\noindent\textcolor{blue}{[序号]主要责任者. 电子文献题名[文献类型/载体类型]. 电子文献的出版或可获得地址(电子文献地址用文字表述), 发表或更新日期/引用日期(任选). }
\mycite{yaoboyuan}


\textcolor{blue}{关于参考文献的未尽事项可参考国家标准《信息与文献参考文献著录规则》（GB/T 7714—2015）}

%输出所有的参考文献
%\printbibliography[heading=none]
注释掉即可，这个只是告诉你参考文献的类别以及规范
	% 参考文献

% 参考文献开始
\unnumchapter{参考文献使用规则}
\renewcommand{\thechapter}{参考文献}

% 设置参考文献字号为 5 号
\renewcommand*{\bibfont}{\zihao{5}}
% 设置参考文献各个项目之间的垂直距离为 0
\setlength{\bibitemsep}{0ex}
\setlength{\bibnamesep}{0ex}
\setlength{\bibinitsep}{0ex}
% 设置单倍行距
\renewcommand{\baselinestretch}{1.2}
% 设置参考文献顺序标签 `[1]` 与文献内容 `作者. 文献标题...` 的间距
\setlength{\biblabelsep}{0.5mm}
% 设置参考文献后文缩进为 0（与 Word 模板保持一致）
\renewcommand{\itemcmd}{
  \addvspace{\bibitemsep} % 恢复 \bibitemsep 的作用
  \mkgbnumlabel{\printfield{labelnumber}}
  \hspace{\biblabelsep}}
% 删除默认的「参考文献 / Reference」标题，使用上面定义的 section 标题

\textcolor{blue}{参考文献书写规范}

\textcolor{blue}{参考国家标准《信息与文献参考文献著录规则》【GB/T 7714—2015】，参考文献书写规范如下：}

\textcolor{blue}{\textbf{1. 文献类型和标识代码}}

\textcolor{blue}{普通图书：M}\qquad\textcolor{blue}{会议录：C}\qquad\textcolor{blue}{汇编：G}\qquad\textcolor{blue}{报纸：N}

\textcolor{blue}{期刊：J}\qquad\textcolor{blue}{学位论文：D}\qquad\textcolor{blue}{报告：R}\qquad\textcolor{blue}{标准：S}

\textcolor{blue}{专利：P}\qquad\textcolor{blue}{数据库：DB}\qquad\textcolor{blue}{计算机程序：CP}\qquad\textcolor{blue}{电子公告：EB}

\textcolor{blue}{档案：A}\qquad\textcolor{blue}{舆图：CM}\qquad\textcolor{blue}{数据集：DS}\qquad\textcolor{blue}{其他：Z}

\textcolor{blue}{\textbf{2. 不同类别文献书写规范要求}}

\textcolor{blue}{\textbf{期刊}}

\noindent\textcolor{blue}{[序号]主要责任者. 文献题名[J]. 刊名, 出版年份, 卷号(期号): 起止页码. }


\textcolor{blue}{\textbf{普通图书}}

\noindent\textcolor{blue}{[序号]主要责任者. 文献题名[M]. 出版地: 出版者, 出版年. 起止页码. }
\mycite{Raymer1992Aircraft}


\textcolor{blue}{\textbf{会议论文集}}

\noindent\textcolor{blue}{[序号]析出责任者. 析出题名[A]. 见(英文用In): 主编. 论文集名[C]. (供选择项: 会议名, 会址, 开会年)出版地: 出版者, 出版年. 起止页码. }
\mycite{sunpinyi}


\textcolor{blue}{\textbf{专著中析出的文献}}

\noindent\textcolor{blue}{[序号]析出责任者. 析出题名[A]. 见(英文用In): 专著责任者. 书名[M]. 出版地: 出版者, 出版年.起止页码. }
\mycite{luoyun}


\textcolor{blue}{\textbf{学位论文}}

\noindent\textcolor{blue}{[序号]主要责任者. 文献题名[D]. 保存地: 保存单位, 年份. }
\mycite{zhanghesheng,Sobieski}


\textcolor{blue}{\textbf{报告}}

\noindent\textcolor{blue}{[序号]主要责任者. 文献题名[R]. 报告地: 报告会主办单位, 年份. }
\mycite{fengxiqiao,Sobieszczanski}


\textcolor{blue}{\textbf{专利文献}}

\noindent\textcolor{blue}{[序号]专利所有者. 专利题名[P]. 专利国别: 专利号, 发布日期. }
\mycite{jiangxizhou}


\textcolor{blue}{\textbf{国际、国家标准}}

\noindent\textcolor{blue}{[序号]标准代号. 标准名称[S]. 出版地: 出版者, 出版年. }
\mycite{GB/T16159—1996}


\textcolor{blue}{\textbf{报纸文章}}

\noindent\textcolor{blue}{[序号]主要责任者. 文献题名[N]. 报纸名, 出版年, 月(日): 版次. }
\mycite{xiexide}


\textcolor{blue}{\textbf{电子文献}}

\noindent\textcolor{blue}{[序号]主要责任者. 电子文献题名[文献类型/载体类型]. 电子文献的出版或可获得地址(电子文献地址用文字表述), 发表或更新日期/引用日期(任选). }
\mycite{yaoboyuan}


\textcolor{blue}{关于参考文献的未尽事项可参考国家标准《信息与文献参考文献著录规则》（GB/T 7714—2015）}

%输出所有的参考文献
%\printbibliography[heading=none]

	% 参考文献
% 参考文献开始
\unnumchapter{参考文献}
\renewcommand{\thechapter}{参考文献}

% 设置参考文献字号为 5 号
\renewcommand*{\bibfont}{\zihao{5}}
% 设置参考文献各个项目之间的垂直距离为 0
\setlength{\bibitemsep}{0ex}
\setlength{\bibnamesep}{0ex}
\setlength{\bibinitsep}{0ex}
% 设置单倍行距
\renewcommand{\baselinestretch}{1.2}
% 设置参考文献顺序标签 `[1]` 与文献内容 `作者. 文献标题...` 的间距
\setlength{\biblabelsep}{0.5mm}
% 设置参考文献后文缩进为 0（与 Word 模板保持一致）
\renewcommand{\itemcmd}{
  \addvspace{\bibitemsep} % 恢复 \bibitemsep 的作用
  \mkgbnumlabel{\printfield{labelnumber}}
  \hspace{\biblabelsep}}
%从这里开始编写---------------------------------------------------------------------
% 删除默认的「参考文献 / Reference」标题，使用上面定义的 section 标题


%输出所有的参考文献
\printbibliography[heading=none]

	\newpage
	
	% 附录：根据自身论文，修改附录的Tex文件（补充论文，不是必须）
	% 附录
% !TeX root = ../zufe_report.tex
\setlength{\baselineskip}{20pt}
\unnumchapter{附录A~~~~附录内容的名称}
\renewcommand{\thechapter}{附录}

% % 设置附录编号格式
% \ctexset{
%   section/number = 附录\Alph{section}
% }

% 附录相关内容…

% % 这里示范一下添加多个附录的方法：

% \section{\LaTeX 环境的安装}
% \LaTeX 环境的安装。

% \section{使用说明}
% 使用说明。

\textcolor{blue}{以下内容可放在附录之内：
1．正文内过于冗长的公式推导；
2．方便他人阅读所需的辅助性数学工具或表格；
3．重复性数据和图表；
4．论文使用的主要符号的意义和单位；
5．程序说明和程序全文;
6．调研报告。
}

\textcolor{blue}{这部分内容可省略（如果省略，删去此页）。阅后删除此段。}

	\newpage

	% 致谢：根据自身要求，修改致谢的Tex文件
	% 致谢
% !TeX root = ../zufe_report.tex
\setlength{\baselineskip}{20pt}
\unnumchapter{致~~~~谢}
\renewcommand{\thechapter}{致谢}

\ctexset{
  section/number = \arabic{section}
}

% 致谢部分尽量不使用 \subsection 二级标题，只使用 \section 一级标题

值此论文完成之际，首先向我的导师……

\textcolor{blue}{致谢正文样式与文章正文相同：宋体、小四；行距：22 磅；间距段前段后均为 0 行。阅后删除此段。}

\end{document}
